% Source: https://artofproblemsolving.com/community/c3h1365852_inscribed_circles
% Stable identity: aops-topic-1365852-q01
Since $\angle A=\angle B=90^\circ$, the sides $AD$ and $BC$ are parallel and $AB$ is perpendicular to both, so $AB$ is the height. Let the inscribed circle have radius $r$; it touches $AD$ and $BC$, so $AB=2r$.

Tangent segments from a point to a circle are equal. From $D$ both tangents have length $DE=4$, and from $C$ both have length $CE=16$. From $A$ and from $B$ the two tangents have length $r$ (the tangency points and the centre form squares at the right-angled corners). So
\[
AD=r+4,\qquad BC=r+16,\qquad CD=20 .
\]

\begin{asy}
usepackage("amsmath");
size(6.5cm);
defaultpen(fontsize(10pt));
pair A=(0,0), B=(0,16), C=(24,16), D=(12,0), O=(8,8);
pair Ep=D+4/20*(C-D);
draw(A--B--C--D--cycle,linewidth(0.9));
draw(circle(O,8),blue);
draw(O--Ep,red+dashed); draw(O--(0,8),gray+dashed); draw(O--(8,0),gray+dashed);
draw((12,0)--(12,16),gray+dotted);
dot(O); dot(Ep);
label("$A$",A,SW); label("$B$",B,NW); label("$C$",C,NE); label("$D$",D,SE); label("$E$",Ep,E);
label("$16$",(C+Ep)/2,E); label("$4$",(D+Ep)/2,E);
label("$r$",(O+Ep)/2,N,red);
draw((0,1)--(1,1)--(1,0)); draw((0,15)--(1,15)--(1,16));
label("\mbox{}",(min(currentpicture,user=true)+max(currentpicture,user=true))/2,black); // keeps the text after the figure black when LaTeX typesets the labels inline
\end{asy}

Drop a perpendicular from $D$ to $BC$ (dotted). It has length $2r$, and it cuts off $BC-AD=12$ along $BC$, so by the Pythagorean theorem
\[
(2r)^2+12^2=20^2\quad\Longrightarrow\quad 2r=16,\ r=8 .
\]
Then $AD=12$, $BC=24$, $AB=16$ and
\[
[ABCD]=\dfrac{(12+24)\cdot16}{2}=\boxed{288}.
\]
(As a check, $r^2=CE\cdot ED=64$, a known property of a circle inscribed in a right trapezoid.)
